\chapter{电路中的电流定律与电压定律}\label{ch:g8:circuit-laws}

两台仪表，两个量 --- 现在，终于轮到
“宪法”登场了。你将搭建的每一个电路，从手电筒
到城市电网，都遵循四条简短的定律：两条管电流的账，
两条管电压的账。掌握它们，再借助一点代数，
你就能算出尚未测量的读数 --- 也能
找出任何“说谎”的测量结果。

\section{电流定律}

\begin{proposition}[串联回路的电流定律]\label{prop:g8:circuit-laws:seriesI}
沿单一回路，各处的电流强度都相同：
\[
I_{\text{电池}} = I_{\text{灯泡 1}} = I_{\text{灯泡 2}} =
\dots
\]
一条道路，一路前行 --- 上一章的收费站，如今
升格为宪法中的定律。
\end{proposition}

\begin{proposition}[节点定律]\label{prop:g8:circuit-laws:junction}
在任一节点处，同时流入的电流强度之和
等于同时流出的电流强度之和：对于一条干路分成
两条支路的情形，
\[
I = I_1 + I_2 .
\]
电荷既不会在分叉处产生、消失，也不会在此处储存：
来多少，去多少。（对任意多条支路均成立 ---
把流入求和、把流出求和，账目便平衡。）
\end{proposition}

\begin{omfigure}
\begin{tikzpicture}[scale=0.85]
  \draw (0,0) to[battery1] (0,3.2)
    to[ammeter, l={$I = \qty{0.50}{A}$}] (3.0,3.2);
  \draw (3.0,3.2) to[short] (3.0,2.3)
    to[lamp, l={$I_1 = \qty{0.30}{A}$}] (6.2,2.3)
    to[short] (6.2,3.2);
  \draw (3.0,3.2) to[short] (3.0,4.1)
    to[lamp, l={$I_2 = \qty{0.20}{A}$}] (6.2,4.1)
    to[short] (6.2,3.2);
  \draw (6.2,3.2) to[short] (7.8,3.2) to[short] (7.8,0)
    to[short] (0,0);
  \fill (3.0,3.2) circle (2pt);
  \fill (6.2,3.2) circle (2pt);
\end{tikzpicture}

{\small 节点定律的数值示例：\qty{0.50}{A} 流至
分叉处，$0.30 + 0.20$ 沿两条支路流出 --- 而
\qty{0.50}{A} 在远端节点处重新汇合。}
\end{omfigure}

\section{电压定律}

\begin{proposition}[串联回路的电压定律]\label{prop:g8:circuit-laws:seriesU}
在串联回路中，电池的电压在各用电器之间分配，
各部分之和恰好等于总电压：
\[
U_{\text{电池}} = U_1 + U_2 + \dots
\]
瀑布的比喻同样能说明这笔账：电池把
电荷的行进抬升至满高度，各用电器则
逐滴地分掉这段高度 --- 良好的导线
则一分也不占用。
\end{proposition}

\begin{proposition}[并联支路的电压定律]\label{prop:g8:circuit-laws:parallelU}
并联的用电器都跨接在同样的两个节点之间，因此
同一电压对它们均适用：
\[
U_1 = U_2 = U_{\text{两节点之间}} .
\]
每条支路都获得两节点间的全部“推力” ---
这正是并联灯泡都能以额定亮度发光的原因，
也是家中每个插座提供相同
\qty{230}{V}。
\end{proposition}

\begin{omfigure}
\begin{tikzpicture}[scale=0.85]
  \draw (0,0) to[battery1, l=\qty{4.5}{V}] (0,3.0)
    to[short] (1.8,3.0)
    to[lamp, l={$U_1 = \qty{2.1}{V}$}] (4.6,3.0)
    to[lamp, l={$U_2 = \qty{2.4}{V}$}] (7.4,3.0)
    to[short] (8.8,3.0)
    to[short] (8.8,0)
    to[short] (0,0);
\end{tikzpicture}

{\small 串联电压定律的数值示例：电池的
\qty{4.5}{V} 被分成两份不等的部分 --- $2.1 + 2.4 =
4.5$，在小数点后一位上完全吻合。}
\end{omfigure}

\begin{example}[为什么分配不均？]\label{ex:g8:circuit-laws:unequal}
两个不同的灯泡串联时通过的电流相同（定律一），
但电压却分配不均 --- \qty{2.1}{V} 与
\qty{2.4}{V}。较大的那一份，归于对电荷行进阻碍更大的灯泡，
也就是说，阻碍更大的灯泡承担更大的电压降。
所谓“阻碍更大”究竟指什么，以及如何测量它，
正是下一章的全部主题 --- 本章的定律
则是它的舞台。
\end{example}

\section{先算后测}

\begin{method}[用四条定律求解电路]\label{met:g8:circuit-laws:solve}
已知某电路中的部分读数：
\begin{enumerate}
  \item 画出电路图，并标出所有已知的 $I$ 和 $U$；
  \item 将适用的定律写成方程 --- 节点处电流求和、
        回路中电压求和，并写出相同性条件：
        串联时 $I$、并联时 $U$ 相同；
  \item 解出未知量（代数课上学过的
        方程，此时终于派上用场）；
  \item 检验结果是否合理：没有哪个灯泡获得的电压超过电池所能提供的，
        也没有哪条支路的电流超过干路的电流。
\end{enumerate}
\end{method}

\begin{example}[例题：混联电路]\label{ex:g8:circuit-laws:worked}
一节 \qty{6}{V} 的电池给灯泡 A 供电，灯泡 A 与
并联的一对灯泡（B 和 C）串联。测得 $I_{\text{电池}} = \qty{0.5}{A}$，
$I_B = \qty{0.3}{A}$，$U_A = \qty{2.5}{V}$。求出其余的量。
节点电流定律：$I_C = 0.5 - 0.3 = \qty{0.2}{A}$。串联电流
定律：灯泡 A 通过全部 \qty{0.5}{A} 的电流。回路电压定律：
这一并联部分两端的电压为 $U_{\text{并联部分}} = 6 - 2.5 = \qty{3.5}{V}$；
并联电压定律：$U_B = U_C = \qty{3.5}{V}$。电路中每只
电表的读数现在都已知——用了四条定律，每条只需一行
代数运算。
\end{example}

\begin{example}[用定律辨别真伪]\label{ex:g8:circuit-laws:liedetector}
某实验报告称：电池电压 \qty{4.5}{V}；两只串联灯泡的电压分别为
\qty{2.1}{V} 和 \qty{2.9}{V}。按回路电压定律把各份电压相加：
共 \qty{5.0}{V}，而电池只“付出”了 \qty{4.5}{V}——不可能；
有人读错了表盘。另一份报告称：节点处电流为
\qty{0.40}{A}，两条支路电流分别为 \qty{0.25}{A} 和 \qty{0.20}{A}
——分叉处“凭空多出”了 \qty{0.05}{A}，予以否定。在定律出现之前，
错误的读数看起来与正确的一样可靠；如今，
这部“宪法”逐页审核每一个数据。
\end{example}

\begin{remark}[同样定律下的电池]\label{rem:g8:circuit-laws:batteries}
“首尾相接”的规则，正是串联电压定律穿着它最古老的
“外衣”：串联电池的电压相加（$1.5 +
1.5 + 1.5 = \qty{4.5}{V}$——这正是这组电池的秘密），
而一节接反的电池会\emph{减去}电压，它的“推力”用于对抗其余电池，
定律对电压的提供者与消耗者一视同仁，区别只在于
它们贡献的正负号不同。
\end{remark}

\section{练习}

\begin{exercise}[$\star$]\label{exo:g8:circuit-laws:1}
写出四条定律——电流的两条、电压的两条——
每条用一句话表述。
\end{exercise}

\begin{exercise}[$\star$]\label{exo:g8:circuit-laws:2}
某节点流入 \qty{0.8}{A}，沿一条支路流出 \qty{0.55}{A}。
求另一条支路的电流。
\end{exercise}

\begin{exercise}[$\star$]\label{exo:g8:circuit-laws:3}
一节 \qty{9}{V} 的电池给三只串联灯泡供电，其中两只分得的电压分别为
\qty{2.5}{V} 和 \qty{3.8}{V}，第三只灯泡的电压是多少？
\end{exercise}

\begin{exercise}[$\star$]\label{exo:g8:circuit-laws:4}
两只灯泡并联后接在 \qty{12}{V} 的电池两端：每只灯泡
两端的电压各是多少？家用插座的哪一个日常事实
遵循的是同样的定律？
\end{exercise}

\begin{exercise}[$\star$]\label{exo:g8:circuit-laws:5}
在前面的混联电路例题中，哪一个再补测的量会是多余的？
为什么？（任选一个量，
并用一条定律加以说明。）
\end{exercise}

\begin{exercise}[$\star$]\label{exo:g8:circuit-laws:6}
两只不同的串联灯泡通过相同的电流，但分得的电压
分别为 $1.8$ 与 \qty{2.7}{V}。哪只灯泡更“阻碍”电流
前进？下一章将用什么量来度量这种阻碍？
\end{exercise}

\begin{exercise}[$\star$]\label{exo:g8:circuit-laws:7}
核查：电池电压 \qty{6}{V}；串联灯泡的测量值分别为 \qty{2.2}{V}、
\qty{1.9}{V}、\qty{1.4}{V}。在电表准确度范围内，这些
数据是否自洽？
\end{exercise}

\begin{exercise}[$\star$]\label{exo:g8:circuit-laws:8}
核查：干路电流 \qty{0.9}{A}；三条支路的电流分别为 \qty{0.4}{A}、
\qty{0.3}{A}、\qty{0.3}{A}。结论如何？
\end{exercise}

\begin{exercise}[$\star\star$]\label{exo:g8:circuit-laws:9}
一节 \qty{4.5}{V} 的电池，灯泡 D 与一个开关串联——
然后，与灯泡 D 并联一只电压表，它在\emph{断开的}开关两端的读数为 \qty{4.5}{V}，
在 D 两端的读数为 \qty{0}{V}。证明这两个读数都满足
回路电压定律。
\end{exercise}

\begin{exercise}[$\star\star$]\label{exo:g8:circuit-laws:10}
三节相同的电池首尾相接给一只灯泡供电；第四节
反向接入这一串电池。按带符号的串联电压定律，求接入前后
的总电压各为多少？灯泡会怎样？
\end{exercise}

\begin{exercise}[$\star\star$]\label{exo:g8:circuit-laws:11}
一个双支路电路：第一条支路中灯泡 E 与 F 串联，
第二条支路仅有灯泡 G。电池电压 \qty{6}{V}；
$U_E = \qty{2.2}{V}$。求 $U_F$ 与 $U_G$，每一步都要
注明所依据的定律。
\end{exercise}

\begin{exercise}[$\star\star\star$]\label{exo:g8:circuit-laws:12}
完整核查：一节 \qty{6}{V} 的电池；灯泡 P 与
一个节点串联；支路 Q 与 R 再汇合后流回。
测得：$I_P = \qty{0.6}{A}$，$I_Q = \qty{0.45}{A}$，
$U_Q = \qty{2.3}{V}$。求电路中其余的每个电流和电压
（$I_R$、$U_P$、$U_R$），并写出每一步所依据的定律——
然后编造一个粗心的学生可能会报告的额外“测量数据”，
这一数据会使你完成的解答
立刻将其判为错误。
\end{exercise}

\section{周末问题：密封的神秘盒子}

\begin{problem}[{周末问题 --- 电子俱乐部的
密封盒子挑战：四条定律对付四个谜团；
总审计}]\label{pb:g8:circuit-laws:1}
俱乐部把简单电路密封进盒中，只留出电表插孔。
各队须手持“章程”推断盒内电路 ---
工作台电池全程为 \qty{6}{V}。

\medskip\noindent\textbf{第一部分 --- 盒子 A：两个接线柱、两
盏灯。}
标签注明：盒内只有两盏相同的灯，
或串联，或并联。
\begin{enumerate}
  \item 盒子 A 接到电池上时电流为 \qty{0.2}{A}；已知单独一盏同样的灯
        接在 \qty{6}{V} 上时，电流约为 \qty{0.4}{A}。
        串联还是并联？请
        用一条电流定律加以论证。
  \item 预测若改用\emph{另一种}
        接法，盒子将抽取多大电流。
  \item 在实际接法中，每盏内部灯泡分得的电压是多少，
        依据哪一条电压定律？
  \item 在哪种接法下，内部的灯泡会以
        额定亮度发光 --- 为何从盒子的窥孔看去
        它们明显并不亮？
\end{enumerate}

\medskip\noindent\textbf{第二部分 --- 盒子 B：三灯
谜题。}
盒子 B 内为：一盏灯（X）与一对并联的灯
（Y、Z --- 不一定相同）串联。
\begin{enumerate}[resume]
  \item 小组测得 $I_{\text{电池}} = \qty{0.5}{A}$
        并在 Y 自身的测试插孔处测得 $I_Y = \qty{0.35}{A}$。
        求 $I_Z$ 与 $I_X$，并注明所依据的定律。
  \item 将电压表并联在 X 两端，示数为 \qty{2.8}{V}。求并联部分
        的电压 --- 以及 $U_Y$、$U_Z$ 各自的值。
  \item Y 与 Z 在相同电压下通过的电流不同。
        这说明两盏灯有何不同？请用
        \cref{ex:g8:circuit-laws:unequal}中的说法作答？
  \item 小组从舱口拧下灯泡 Y。预测新的
        $I_{\text{电池}}$ 的变化方向（上升、下降
        或不变），并依据“道路”及其“阻碍”进行推理。
\end{enumerate}

\medskip\noindent\textbf{第三部分 --- 盒子 C：作弊。}
\begin{enumerate}[resume]
  \item 盒子 C 的搭建者报告：“电池 \qty{6}{V}；盒内两盏串联的灯，
        电压分别为 \qty{2.5}{V} 与 \qty{4.1}{V}；
        盒内一个节点流入 \qty{0.3}{A}，流出
        \qty{0.2}{A} 与 \qty{0.15}{A}。”请判定该报告有误
        两次，每次依据一条定律。
  \item 两个报告的电压读数之一，后来被证实属实。
        若属实者为 \qty{2.5}{V}，则另一读数
        实际必须是多少？
  \item 搭建者承认节点的数值是
        “慷慨取整”的。若两条支路的读数
        \qty{0.2}{A} 与 \qty{0.15}{A} 属实，则
        干路实际必须通过多大电流？
  \item 四条定律为何使密封盒子里的作弊
        成为一场必输的游戏？用一句话作答。
\end{enumerate}

\medskip\noindent\textbf{第四部分 --- 章程。}
\begin{enumerate}[resume]
  \item 撰写俱乐部章程：用四行写出四条定律，
        每条配以承载它的日常意象（道路、
        收费站、瀑布、岔路口）。
\end{enumerate}
\end{problem}
